% SPDX-License-Identifier: CC-BY-4.0
% Copyright (c) 2026 Aravind Karthigeyan
\documentclass[11pt]{article}
\usepackage[T1]{fontenc}
\usepackage[margin=1in]{geometry}
\usepackage{amsmath,amssymb,amsthm,booktabs,array,longtable}
\usepackage{xurl}
\usepackage[colorlinks=true,linkcolor=blue,citecolor=blue,urlcolor=blue]{hyperref}
\usepackage{microtype}
\setlength{\emergencystretch}{1em}
\newtheorem{theorem}{Theorem}[section]
\newtheorem{lemma}[theorem]{Lemma}
\newtheorem{proposition}[theorem]{Proposition}
\newtheorem{corollary}[theorem]{Corollary}
\theoremstyle{remark}
\newtheorem{remark}[theorem]{Remark}
\newcommand{\R}{\mathbb R}
\newcommand{\Q}{\mathbb Q}
\newcommand{\ip}[2]{\langle #1,#2\rangle}
\newcommand{\norm}[1]{\lVert #1\rVert}
\newcommand{\code}[1]{{\small\nolinkurl{#1}}}
\newcommand{\pathcode}[1]{\nolinkurl{#1}}
\DeclareMathOperator{\tr}{tr}
\title{Global Optimality and Uniqueness for Nine Coulomb Charges on the Sphere:\\
An Exact-Certificate Proof in Lean 4}
\author{Aravind Karthigeyan\\\small University of Illinois Urbana-Champaign}
\date{October 6, 2026\\\small Review draft; not peer reviewed}
\begin{document}
\maketitle
\begin{abstract}
For nine distinct points on the unit sphere in real three-space, the sum of
reciprocal Euclidean distances is minimized by a particular tricapped trigonal
prism. We describe a computer-assisted proof of global optimality and uniqueness
up to orthogonal transformations and relabeling, together with its Lean 4
formalization. The candidate is defined by a radical stationarity equation and
a rational isolating interval, not by numerical coordinates. A rational
degree-20 three-point semidefinite certificate is nearly sharp, with its entire
residual paid explicitly. Exact Bernstein certificates force pair inner products
into narrow windows. Finite Gram-matrix classification gives one coherent
relabeling, and a certified ambient Hessian bound completes the argument locally.
The resulting Lean theorem has no certificate hypotheses. Its recorded axiom
closure consists only of propositional extensionality, classical choice, and
quotient soundness. A separately implemented kernel also accepted the exported
complete theorem and its dependencies. The source development and this draft
received substantial AI assistance. Independent human review and reproduction
by another operator remain requested; no priority claim is made here.
\end{abstract}

\section{Statement and scope}
Let
\[
 \mathcal C_9=\{(x_0,\ldots,x_8)\in(\R^3)^9:
          \norm{x_i}=1\text{ for all }i,\quad x_i\ne x_j\text{ for }i\ne j\}.
\]
The energy convention throughout is
\begin{equation}\label{eq:energy}
 E(X)=\sum_{0\le i<j<9}\frac1{\norm{x_i-x_j}}
     =\sum_{i<j} f(\ip{x_i}{x_j}),\qquad
 f(t)=\frac1{\sqrt{2-2t}}.
\end{equation}
Thus there are 36 terms; neither ordered pairs nor reciprocal squared distances
are used. Denote by $P=P(r_*)$ the configuration defined in Section~\ref{sec:candidate}.

\begin{theorem}\label{thm:main}
For every $X\in\mathcal C_9$, $E(X)\ge E(P)$. Moreover,
\[
 E(X)=E(P)\quad\Longleftrightarrow\quad
 \exists Q\in O(3),\ \exists\sigma\in S_9,
       \forall i,\quad x_i=Qp_{\sigma(i)}.
\]
The parameter $r_*$ exists and is uniquely determined by the exact
characterization below, and $P\in\mathcal C_9$.
\end{theorem}

The formal statement is independent of the global proof: the root file
\pathcode{lean/Challenge.lean} imports the candidate definitions, not the
certificate assembly. It asks for parameter existence and uniqueness, global
optimality, and both directions of the equality characterization. The final
theorem is
\begin{quote}
\code{ThomsonN9.thomson_nine_complete :}\par
\code{ThomsonN9.Challenge.CompleteTarget}.
\end{quote}
The configuration space uses Mathlib's \code{EuclideanSpace} indexed by
\code{Fin 3}, with $\R$ as the scalar field. Congruence uses a real
linear isometry equivalence and a permutation. In particular, reflections
are allowed. Separate fixed-candidate theorems remove the parameter argument.

This paper follows the \emph{selected formal route}: the reduced D60
certificate, Bernstein positivity, finite Gram-window classification, and the
ambient augmented Hessian. Earlier Hermite-remainder and higher-precision
certificate routes remain in the research archive, but are not substituted
for the proof checked here. Large finite data are part of the accompanying
artifact, not printed as thousands of coefficients in this paper. We state
their mathematical obligations and identify their checkers explicitly.
In the publication repository, all proof-source, certificate, and verification
paths below are relative to the \pathcode{proof/} directory unless explicitly
identified as publication-level paths. The supplemental upstream challenge
bridge is instead in the publication-level \pathcode{bridge/} directory.

\subsection{Relationship to prior work}
Cohn and Kumar's universal-optimality theorem gives a conceptual framework for
several classical spherical energy minima, including the tetrahedron,
octahedron, and icosahedron~\cite{cohnkumar}. Schwartz's five-electron proof
combines global interval elimination with local analysis~\cite{schwartz}.
For eight points, Kryvonos, Liehr, and Taylor establish Coulomb optimality and
uniqueness of an optimized square antiprism using exact certificates, with a
Lean formalization~\cite{klt,n8code}. Tran's seven-point development uses
three-point bounds, geometric localization, and local rigidity~\cite{tran}.
These works motivated the present architecture, especially the use of a
nearly sharp global certificate followed by a quantitative local argument.

The N9 implementation does not import the seven- or eight-point optimality
theorems as premises. An Apache-2.0-licensed Schoenberg dependency closure from
the \code{delsarte} project is vendored with attribution and compatibility
changes recorded~\cite{delsarte}. The inspected upstream commits are specified
in the bibliography and artifact. The separate full N8 project rebuilds were
not completed in this investigation; source inspection is not presented as
successful reproduction of those projects.

\section{The exact candidate}\label{sec:candidate}
Put
\[
 e_0=(1,0),\qquad e_1=(-1/2,\sqrt3/2),\qquad
 e_2=(-1/2,-\sqrt3/2).
\]
For $0<r<1$, let $z=\sqrt{1-r^2}$ and define
\begin{equation}\label{eq:candidate}
 p_{k,+}=(re_k,z),\quad p_{k,-}=(re_k,-z),\quad
 c_k=(-e_k,0),\qquad k=0,1,2.
\end{equation}
These are nine distinct unit vectors. Rotations through $2\pi/3$, reflection
in the equatorial plane, and a vertical reflection preserve the configuration.
They generate $D_{3h}$. Under this group the six prism points form one orbit
and the three equatorial caps form the other. The cyclic rotation subgroup
instead has three orbits of size three. Symmetry is a property of this
candidate, not an assumption on a competing configuration.

The six pair types give the following exact spectrum.
\begin{center}
\begin{tabular}{lcc}
\toprule
Pair type & Inner product & Multiplicity\\
\midrule
Same upper or lower triangle & $1-3r^2/2$ & 6\\
Vertically opposite prism points & $2r^2-1$ & 3\\
Other upper/lower pairs & $r^2/2-1$ & 6\\
Two caps & $-1/2$ & 3\\
Cap and opposite prism direction & $-r$ & 6\\
Other cap/prism pairs & $r/2$ & 12\\
\bottomrule
\end{tabular}
\end{center}
Consequently,
\begin{equation}\label{eq:familyenergy}
 F(r)=E(P(r))=
 \frac{2\sqrt3}{r}+\frac3{2\sqrt{1-r^2}}+
 \frac6{\sqrt{4-r^2}}+\sqrt3+
 \frac6{\sqrt{2+2r}}+\frac{12}{\sqrt{2-r}}.
\end{equation}
Define
\begin{align}\label{eq:stationarity}
 S(r)=F'(r)
 &= -\frac{2\sqrt3}{r^2}+\frac{3r}{2(1-r^2)^{3/2}}
    +\frac{6r}{(4-r^2)^{3/2}}\notag\\
 &\hspace{15mm}-\frac6{(2+2r)^{3/2}}+\frac6{(2-r)^{3/2}}.
\end{align}
All square roots here are positive. Let
\[
 A=7105483340929538113277041792727912437680,\qquad
 I=[A/10^{40},(A+1)/10^{40}].
\]
The exact definition of $r_*$ is the unique $r\in I$ with $S(r)=0$.

\begin{lemma}\label{lem:parameter}
There is exactly one zero $r_*$ of $S$ in $I$. It is in the interior of $I$
and is the unique minimizer of $F$ on $(0,1)$.
\end{lemma}
\begin{proof}
Continuity holds on $(0,1)$. Direct differentiation gives
\[
 S'(r)=\frac{4\sqrt3}{r^3}
 +\frac{3(1+2r^2)}{2(1-r^2)^{5/2}}
 +\frac{12(2+r^2)}{(4-r^2)^{5/2}}
 +\frac{18}{(2+2r)^{5/2}}+\frac9{(2-r)^{5/2}}>0.
\]
Exact rational interval evaluations establish $S(A/10^{40})<0$ and
$S((A+1)/10^{40})>0$. These two finite evaluations, including all denominator
signs and square-root enclosures, are checked in the parameter certificate
modules. The intermediate value theorem and strict monotonicity give the
claim. Positivity of $S'$ also gives strict decrease of $F$ before the root
and strict increase afterwards. See \code{parameter_exists_unique} and
\pathcode{FamilyMinimum.lean} for the checked argument.
\end{proof}

The original radical equation is retained instead of replacing it by an
unqualified polynomial equation obtained by squaring. The artifact contains
a separate degree-104 integer elimination polynomial, but no irreducibility
claim is needed and it is not the defining premise of the formal theorem.
For orientation only,
\begin{align*}
 r_*&\approx0.7105483340929538113277041792727912437680,\\
 E(P)&\approx25.7599865312698315640757022737604940.
\end{align*}
Neither decimal is a proof input.

\section{Global scalar certificates}\label{sec:scalar}
For the rest of the proof set
\[
 M=\frac{5593}{10000},\qquad \varepsilon=10^{-4},\qquad\delta=10^{-8}.
\]
An exact rational upper bound $U$ for $E(P)$ is defined in
\pathcode{Certificates/ReducedEnergy.lean}:
\[
 U=\frac{
 5151997306253966312815140454752098805728382727295351889952429}
 {200000000000000000000000000000000000000000000000000000000000}.
\]
It comes from interval evaluation of $F$ at a rational radius and
Lemma~\ref{lem:parameter}. Thus $E(P)\le U$ uses only minimality within the
symmetric family, not Theorem~\ref{thm:main}.

\subsection{Separation by a two-point bound}
\begin{lemma}[Separation]\label{lem:separation}
If $X\in\mathcal C_9$ and $E(X)\le E(P)$, then
$\ip{x_i}{x_j}<M$ for all $i\ne j$.
\end{lemma}
\begin{proof}[Certificate proof]
The separation polynomial $q$ has a Legendre expansion
$q(t)=\sum_k a_k P_k(t)$, with $P_k(1)=1$. Its nonzero coefficients are
\[
\begin{array}{c|rrrrr}
k&0&1&2&3&7\\\hline
10^{12}a_k&806045161636&460032960362&200888085899&63087598053&5677717462
\end{array}
\]
\[
\begin{array}{c|rrrr}
k&12&17&22&27\\\hline
10^{12}a_k&490951549&542566025&88576655&134439481
\end{array}
\]
The scalar certificate proves $q(t)\le f(t)$ on $[-1,1)$.
Schoenberg positivity gives
\[
 \sum_{i<j}q(\ip{x_i}{x_j})\ge B,
 \qquad B=\frac{81a_0-9q(1)}2.
\]
The separate tail certificate proves, for $M\le t<1$,
$f(t)-q(t)\ge T-B+10^{-8}$, where
\[
 T=\frac{6439996633}{250000000},\qquad
 B=\frac{25728382789209}{10^{12}},\qquad E(P)\le T.
\]
If one pair were in this tail, summing its excess together with the
nonnegative excesses of all other pairs would yield $E(X)>T$, a contradiction.
The literal polynomial/Legendre linkage and both positivity certificates are
checked in \pathcode{LP/ConcreteLinkage.lean} and
\pathcode{ClaudeL2b/LPAdapterCodex.lean}; the application is
\code{sublevel_separation}.
\end{proof}

\subsection{Three-point positivity and an explicit residual}
For a triple of unit vectors write its three inner products as $u,v,t$.
They belong to the capped Gram domain
\[
 \mathcal D_M=\{(u,v,t)\in[-1,M]^3:
       g(u,v,t):=1+2uvt-u^2-v^2-t^2\ge0\}.
\]
The three-dimensional rooted kernels use the polynomial recurrence
\[
 Q_0=1,\quad Q_1=t-uv,\quad
 Q_{k+2}=2(t-uv)Q_{k+1}-(1-u^2)(1-v^2)Q_k.
\]
Products with univariate features and rational rectangular Gram factors give
positive moment blocks. After symmetrizing over the six edge orders, the
repeated-index correction for nine points is
\[
 \mathcal K=7\mathcal S(u,v,t)+\mathcal S(u,u,1)
       +\mathcal S(v,v,1)+\mathcal S(t,t,1)
       +\mathcal S(1,1,1)/8.
\]
The proved positivity assertion is a nonnegative \emph{sum} of $\mathcal K$
over ordered distinct triples of a configuration. Individual values of
$\mathcal K(u,v,t)$ need not be nonnegative.

The reduced certificate supplies a degree-20 rational polynomial $h$ and an
identity on $\mathcal D_M$ of the form
\begin{equation}\label{eq:sdpidentity}
 \frac{h(u)+h(v)+h(t)}3-\frac U{36}
 =\sum_b\mathcal K_b(u,v,t)+\sum_s w_s(u,v,t)q_s(u,v,t)^2
        +\sum_\alpha c_\alpha u^{\alpha_1}v^{\alpha_2}t^{\alpha_3}.
\end{equation}
The weights are nonnegative rational multiples and permutations of
$1$, $1+u$, $M-u$, $g$, $(1+u)(M-u)$, and $g(1+u)(M-u)$.
The factor matrices are rational, so their positive semidefiniteness follows
from their factored form, without a numerical eigenvalue test.

Let $C$ be the rational value \code{Claude.L1.certifiedEnergyReduced} and
$L=U-C$. The coefficient bound certifies
$36\sum_\alpha|c_\alpha|\le L<10^{-46}$. Its approximate value is
$2.396\cdot10^{-54}$; importantly, it is not zero.

\begin{lemma}[Near-sharp pair bound]\label{lem:pairbound}
For nine unit vectors whose distinct-pair inner products are at most $M$,
\[
 U-L\le\sum_{i<j}h(\ip{x_i}{x_j}).
\]
\end{lemma}
\begin{proof}
Sum~\eqref{eq:sdpidentity} over all $9\cdot8\cdot7=504$ ordered distinct
triples. The left side becomes
\[
 14\left(\sum_{i<j}h(\ip{x_i}{x_j})-U\right).
\]
The kernel sum and weighted squares are nonnegative. Every residual monomial
has absolute value at most one, so the right side is at least
$-504\sum_\alpha|c_\alpha|$. Dividing by 14 gives the claimed bound.
\end{proof}

The coefficient identity is a finite integer verification, not a comparison
at sample points. Its degree box has $21^3=9261$ entries and covers all
$\binom{23}{3}=1771$ monomials of total degree at most 20, including coefficients
that must be zero. The final instantiation is
\code{Claude.L1.reduced_pair_bound}. The D60 label refers to the $10^{60}$
rational rounding scale, not to the polynomial degree.

\subsection{Bernstein minorants and exterior gaps}
Write $s_1(r),\ldots,s_6(r)$ for the six inner products in the candidate
spectrum. The certificate establishes the following exact scalar assertions.
\begin{lemma}[Minorant and gaps]\label{lem:gaps}
The reduced polynomial satisfies $h(t)\le f(t)$ for $-1\le t<1$.
For every $r\in I$:
\begin{enumerate}
\item If $-1+\varepsilon\le t\le M$ and
      $|t-s_j(r)|\ge\varepsilon$ for all $j$, then
      $f(t)-h(t)\ge10^{-46}$.
\item If $-3/4\le t\le M$ and
      $|t-s_j(r)|\ge\delta$ for all $j$, then
      $f(t)-h(t)\ge10^{-33}$.
\end{enumerate}
\end{lemma}
\begin{proof}[Exact finite verification]
For a polynomial $a$ and a gap $\gamma$, substitute
$y=\sqrt{(1-t)/2}$; then $0<y\le1$, $t=1-2y^2$, and $f(t)=1/(2y)$.
It suffices to prove
\[
 1-2y\bigl(a(1-2y^2)+\gamma\bigr)\ge0.
\]
For degree 20 this is a polynomial of degree at most 41. On a rational
subinterval it is transformed affinely to $[0,1]$ and expanded in the
Bernstein basis. Nonnegative Bernstein coefficients imply nonnegativity on
that whole interval. Complete subdivision trees are checked, including their
coverage; accepting a few selected leaves would not suffice.

The implemented route first proves a rational precursor minorant and its
coarse gaps, then proves an exact coefficient inequality transferring them
downwards to $h$. Refined gaps are checked with their own finite data.
Rational enclosures of $s_j(r)$ enlarge the exterior intervals to be checked,
so no unexamined interval between contact windows is discarded.
The checked endpoints are \code{Claude.L2b.hNew_le_inv_dist},
\code{Claude.L2b.exterior_gap_new}, and
\code{Claude.L2b.refined_gap_new}.
\end{proof}

For a configuration in the sublevel $E(X)\le E(P)$, Lemmas~\ref{lem:separation}
and~\ref{lem:pairbound} imply
\begin{equation}\label{eq:slack}
 0\le\sum_{i<j}\bigl(f(\ip{x_i}{x_j})-h(\ip{x_i}{x_j})\bigr)\le L<10^{-46}.
\end{equation}
Each summand is nonnegative. Hence a single exterior pair from the first gap
assertion would contradict~\eqref{eq:slack}. Every pair therefore belongs to a
contact window of radius $\varepsilon$ or the retained interval
$[-1,-1+\varepsilon]$. Exclusion of that last interval is a geometric step,
not an assumption of the minorant certificate.

\section{Finite Gram rigidity and refined localization}\label{sec:gram}
\begin{lemma}[Seven-window rigidity]\label{lem:rigidity}
Let $r\in I$, and let $x_0,\ldots,x_8$ be unit vectors in $\R^3$. Suppose
every distinct-pair inner product lies within $\varepsilon$ of some $s_j(r)$
or lies in $[-1,-1+\varepsilon]$. Then the latter possibility never occurs.
There is a permutation $\sigma$ such that
\[
 |\ip{x_{\sigma(i)}}{x_{\sigma(j)}}-
             \ip{P(r)_i}{P(r)_j}|\le\varepsilon\qquad\text{for all }i,j.
\]
If the original six contact windows have a common radius
$0\le\eta\le\varepsilon$, the same conclusion holds with $\eta$.
\end{lemma}
\begin{proof}[Finite classification with checked geometric soundness]
The six rationally enclosed windows are pairwise disjoint, so they assign a
unique color to each edge. First, for a near-antipodal pair and a third point,
the two other edges lie in one of $7^2=49$ pairs of windows. Exact interval
evaluation in each case contradicts nonnegativity of the three-point Gram
determinant. Next, the six-window hypothesis forces a triple with Gram
determinant greater than $1/100$. The anchor-free certificate checks the
remaining colored four-point possibilities and uses the vanishing of their
four-dimensional Gram determinant in $\R^3$.

For an invertible anchor Gram matrix $G_A$ and another point with inner-product
column $v$ against the anchor, the unit-norm condition gives
\[
 v^{\mathsf T}\operatorname{adj}(G_A)v=\det G_A.
\]
For two additional points with columns $v,w$ and mutual inner product $t$,
\[
 v^{\mathsf T}\operatorname{adj}(G_A)w=t\det G_A.
\]
Interval versions of these polynomial equalities reject impossible edge-color
extensions. A finite coverage theorem checks all remaining anchor and extension
records and shows that they have exactly the candidate's edge pattern up to
one permutation. This step does not choose 36 independent nearest contacts.
Its exact interval operations, geometric implications, exhaustive coverage,
and reference-color identification are proved separately in the
\pathcode{ClaudeWindow/} modules.

For the narrower-window conclusion, uniqueness of the coarse colors ensures
that the same color assigned by the classification is the one supplied by
the narrow hypothesis. This transfers the smaller error without repeating
the finite enumeration.
\end{proof}

\begin{proposition}[Refined sublevel localization]\label{prop:localization}
If $X\in\mathcal C_9$ and $E(X)\le E(P)$, then there is one permutation
$\sigma$ with
\begin{equation}\label{eq:gramclose}
 |\ip{x_{\sigma(i)}}{x_{\sigma(j)}}-\ip{p_i}{p_j}|\le10^{-8}
 \quad\text{for all }i,j.
\end{equation}
\end{proposition}
\begin{proof}
Apply Lemma~\ref{lem:rigidity} to the coarse localization obtained from
\eqref{eq:slack}. Its coherent pattern implies every pair inner product is at
least $-3/4$: the smallest candidate contact is $r_*^2/2-1$, and its coarse
window is above that bound. The second assertion of Lemma~\ref{lem:gaps} now
applies. Since $L<10^{-46}<10^{-33}$, no pair can lie outside the windows of
radius $\delta$. The narrow version of Lemma~\ref{lem:rigidity} gives the
conclusion. The formal result is \code{sublevel_refined_gram}.
\end{proof}

\section{Alignment and the local energy gap}\label{sec:local}
\begin{lemma}[Quantitative orthogonal alignment]\label{lem:alignment}
Under~\eqref{eq:gramclose}, there exist $Q\in O(3)$ and the corresponding
relabeling such that, with $y_i=Qx_{\sigma(i)}$, $b_i=y_i-p_i$, and
$\rho^2=\sum_i\norm{b_i}^2$, one has
\begin{equation}\label{eq:alignment}
 \rho\le10^{-4},\qquad
 \sum_i\bigl(b_i[a]p_i[c]-b_i[c]p_i[a]\bigr)=0
       \quad(a,c\in\{0,1,2\}).
\end{equation}
\end{lemma}
\begin{proof}[Structure of the checked estimate]
The first three prism points give a nondegenerate reference frame. Constructing
orthonormal frames and propagating entrywise Gram errors bounds the best
orthogonal displacement by $10000\delta=10^{-4}$. The displacement minimum
over $O(3)$ exists by compactness. Varying a minimizing transformation by
planar rotations gives the second identity in~\eqref{eq:alignment}. Thus the
rotation gauge is obtained for an arbitrary competitor, rather than imposed
as an extra hypothesis on its coordinates. The quantitative frame inequalities
are in \pathcode{GramAlignment.lean} and \pathcode{CandidateAlignment.lean};
the minimizing-transformation argument is in \pathcode{Alignment.lean}.
\end{proof}

The local calculation is performed on 27 ambient coordinates, avoiding a
numerically chosen basis for the 15 nonrotational tangent directions. For
$d_{ij}=p_i-p_j$ put
\[
 B(d)=\frac{3dd^{\mathsf T}}{\norm d^5}-\frac{I_3}{\norm d^3},\qquad
 \lambda_i=-\frac12\sum_{j\ne i}\frac1{\norm{d_{ij}}}.
\]
Let $H$ be the block Laplacian of the $B(d_{ij})$, minus the block diagonal
matrix with blocks $\lambda_iI_3$. It is the constrained Hessian, with the
radial Lagrange-multiplier correction included. Let $R$ be the $27\times3$
matrix of infinitesimal rotation columns at $P$. The gauge in
\eqref{eq:alignment} is precisely $R^{\mathsf T}b=0$.

\begin{lemma}[Certified augmented Hessian]\label{lem:hessian}
For the exact candidate parameter,
\begin{equation}\label{eq:psd}
 H+RR^{\mathsf T}-\frac7{50}I_{27}\succeq0.
\end{equation}
\end{lemma}
\begin{proof}[Exact interval and congruence certificate]
An arithmetic circuit evaluates enclosures for the entries of the matrix in
\eqref{eq:psd}, uniformly for the entire rational interval $I$, using the
integer scale $2^{180}$. The circuit contains 37,815 checked nodes. Separate
semantic lemmas identify the circuit entries with the actual coordinates,
distances, radial coefficients, pair Hessians, and rotation columns.

For the resulting symmetric matrix $A$, a rational unit upper triangular
matrix $C$ is supplied. The certificate evaluates $C^{\mathsf T}AC$ through
1,458 dot products and 729 final enclosures. Its diagonal entries are at least
$1/1000$ and its off-diagonal entries have absolute value at most $1/100000$.
There are 26 off-diagonal entries in each row, so the diagonal-dominance
margin is
\[
 \frac1{1000}-\frac{26}{100000}=\frac{37}{50000}>0.
\]
Strict diagonal dominance proves positivity; invertibility of the unit upper
triangular $C$ transfers it to $A$. Denominator positivity and directions of
all integer-to-rational enclosure comparisons are proved in the checking
library. The final theorem is
\code{Certificates.AmbientKernel.candidate_augmented_hessian_posSemidef}.
No floating-point spectrum or condition-number estimate is a premise.
\end{proof}

Full candidate stationarity is proved from~\eqref{eq:stationarity} in
\pathcode{LocalStationarity.lean}; it is not inferred just from a numerical
gradient. In the above notation the energy gradient at $p_i$ is
$\lambda_ip_i$. Also every candidate pair distance is at least one.

\begin{lemma}[Local coercivity]\label{lem:coercivity}
Suppose $Y\in\mathcal C_9$, $b_i=y_i-p_i$, $\rho^2=\sum_i\norm{b_i}^2$,
$0\le\rho\le1/500$, and $R^{\mathsf T}b=0$. Then
\begin{equation}\label{eq:localgap}
 E(Y)-E(P)\ge\frac{\rho^2}{50}.
\end{equation}
\end{lemma}
\begin{proof}
An algebraic reciprocal-distance estimate, valid for $\norm d\ge1$ and
$\norm w\le\norm d/200$, is
\[
 \frac1{\norm{d+w}}\ge\frac1{\norm{d}}-\frac{\ip{d}{w}}{\norm{d}^3}
 +\frac{3(\ip{d}{w})^2-\norm{d}^2\norm{w}^2}{2\norm{d}^5}
 -\frac{51}{50}\norm{w}^3.
\]
Apply this to $d=d_{ij}$ and $w=b_i-b_j$. The sphere identities
$2\ip{p_i}{b_i}=-\norm{b_i}^2$ and stationarity convert the linear sum to
the radial correction in $H$. Moreover,
$\sum_{i<j}\norm{b_i-b_j}^3\le(27/2)\rho^3$.
Thus the formal algebraic estimate is
\[
 E(Y)-E(P)\ge\frac12b^{\mathsf T}Hb-\frac{1377}{100}\rho^3.
\]
By Lemma~\ref{lem:hessian} and $R^{\mathsf T}b=0$,
$b^{\mathsf T}Hb\ge(7/50)\rho^2$. Finally,
$7/100-(1377/100)\rho\ge1/50$ for $0\le\rho\le1/500$.
This proves~\eqref{eq:localgap}. All analytic-to-algebraic transitions are
proved in \pathcode{LocalPotential.lean}, \pathcode{LocalEstimate.lean}, and
\pathcode{LocalCoercivity.lean}; no numerical Taylor remainder is used.
\end{proof}

\section{Global conclusion and stability}
\begin{proof}[Proof of Theorem~\ref{thm:main}]
Consider any $X\in\mathcal C_9$. If $E(X)>E(P)$, the lower bound and the
failure of equality are immediate. Otherwise Proposition~\ref{prop:localization}
and Lemma~\ref{lem:alignment} give a congruent configuration $Y$ satisfying
the gauge and $\rho\le10^{-4}<1/500$. Congruence preserves energy.
Lemma~\ref{lem:coercivity} now gives
\[
 0\ge E(X)-E(P)=E(Y)-E(P)\ge\rho^2/50\ge0.
\]
Hence $\rho=0$, $Y=P$, and $X$ is congruent to $P$. Conversely, any orthogonal
image and relabeling of $P$ has its energy. Parameter existence and
admissibility were established independently in Section~\ref{sec:candidate}.
This completes all conjuncts of the formal target without an unproved
localization, positivity, or symmetry hypothesis.
\end{proof}

The available certificate slack also gives a quantitative statement. Define
the squared distance to the candidate orbit by
\[
 d_{\mathrm{orb}}(X,P)^2=
 \min_{Q\in O(3),\sigma\in S_9}\sum_i\norm{Qx_{\sigma(i)}-p_i}^2.
\]
\begin{corollary}[Quantitative stability]\label{cor:stability}
If $X\in\mathcal C_9$ and $E(X)\le E(P)+10^{-47}$, then
\[
 d_{\mathrm{orb}}(X,P)^2\le50\bigl(E(X)-E(P)\bigr).
\]
\end{corollary}
The checked proof repeats the slack localization with the added $10^{-47}$
budget and applies the same local bound. Its public theorem is
\code{thomson_nine_quantitative_stability}. This corollary was checked in
Lean, but the recorded separate-kernel export described next selected the
complete optimality/uniqueness theorem, not this additional root.

\section{Formalization, reproduction, and trust boundaries}\label{sec:verification}
Paths in the following table are under \pathcode{lean/ThomsonN9/}, except the
explicit root and \pathcode{Claude*/} entries. Names in the table are theorem
or module identifiers, not additional mathematical assumptions.
\begin{center}\small
\begin{tabular}{p{.29\linewidth}p{.64\linewidth}}
\toprule
Mathematical component & Formal implementation\\
\midrule
Definition and target & \pathcode{Basic.lean}, \pathcode{Candidate.lean}, root \pathcode{Challenge.lean}\\
Parameter and family & \pathcode{ParameterExists.lean}, \pathcode{CandidateEnergy.lean}, \pathcode{FamilyMinimum.lean}\\
Separation and SDP & \pathcode{LP/Concrete.lean}, \pathcode{SDP/CertificateBound.lean}, \pathcode{ClaudeL1/Main.lean}\\
Bernstein gaps & \pathcode{Bernstein/}, \pathcode{ClaudeL2b/ReducedMinorant.lean}, \pathcode{ClaudeL2b/RefinedWindows.lean}\\
Finite Gram classification & \pathcode{ClaudeWindow/Main.lean}, \pathcode{ClaudeWindow/Bridge.lean}\\
Refined localization & \pathcode{ScalarCertificates.lean}, \pathcode{GlobalReduction.lean}\\
Coordinates and gauge & \pathcode{GramAlignment.lean}, \pathcode{CandidateAlignment.lean}, \pathcode{Alignment.lean}\\
Actual Hessian and PSD & \pathcode{BlockHessian.lean}, \pathcode{Certificates/AmbientKernelAssembly.lean}, \pathcode{Certificates/AmbientHessian.lean}\\
Local and global assembly & \pathcode{LocalEstimate.lean}, \pathcode{LocalCoercivity.lean}, \pathcode{Endgame.lean}, \pathcode{Completion.lean}\\
Additional stability & \pathcode{QuantitativeStability.lean}\\
\bottomrule
\end{tabular}
\end{center}

\subsection{Supplemental bridge to the N7 project's N9 target}
The publication-level \pathcode{bridge/UpstreamN7Bridge.lean} imports the
actual definitions from the pinned N7 challenge and proves its N9 stretch
target, rather than a paraphrase of that target in new definitions. The
translation uses $z=\sqrt{1-r_*^2}$ and the explicit cap relabeling that sends
upstream indices $(6,7,8)$ to candidate indices $(8,6,7)$. An explicit Lean
audit confirmed that the new declaration
\code{ThomsonN9.UpstreamN7Bridge.thomson_nine} has a type definitionally equal
to the upstream N9 target. The isolated source compilation and audit completed
with exit zero on October 6, 2026; the three reported axiom closures contain
only the same standard axioms listed below.

The imported upstream challenge contains unproved placeholders, which produce
expected warnings when that file is compiled in isolation. None is invoked
by the bridge, and no placeholder axiom appears in its proof dependency
closure. The upstream source is obtained separately, not redistributed here.
This is a supplemental Lean check: it is not a run of the separate Lean
Comparator application, and the bridge was not included in the recorded
nanoda export of the frozen N9 complete theorem. The publication-level
\pathcode{bridge/README.md} and \pathcode{bridge/check-result.json} record the
pins, source hashes, scope, and reproduction commands.

\subsection{Finite artifacts}
The exact data are inspectable independently of the numerical generators.
The reduced polynomial and factored identity are in
\pathcode{near_sharp/reduced/near-sharp-d10-reduced-D60.json}; their literal
Lean data and integer checks are in \pathcode{lean/ClaudeL1/}. Its SHA-256 is
\begin{quote}\footnotesize
\pathcode{9816ae66f3b68caf4a03d6c6f7c514124c8cd8ba4415497a79443ce64f158781}.
\end{quote}
The corresponding reduced Bernstein and refined-window data are
\pathcode{near_sharp/reduced/bernstein-reduced.json} and
\pathcode{near_sharp/reduced/refined-windows-reduced.json}.
Finite Gram tables are literal in \pathcode{lean/ClaudeWindow/WindowData.lean}.
The ambient local certificate is accompanied by
\pathcode{certificates/ambient-dd-k0.json}, while the selected Lean endpoint
uses the literal interval circuit and congruence data in the
\pathcode{Certificates/Ambient*} modules. The source manifests bind all data;
the JSON file is not an unchecked oracle imported at proof time.

Numerical SDP solvers, unrestricted numerical minimization, and certificate
generation are discovery tools only. The separate exact Python replay command
\begin{quote}\small
\verb|python verify_global.py --reduced --refined|
\end{quote}
recomputes finite arithmetic premises. It is useful for independent checking,
but a successful Python run is not equivalent to building the Lean proof.
The authoritative, pinned Lean build and audit commands are in
\pathcode{REPRODUCE.md}; final theorem statements, axiom reports, and rejected
negative controls must be checked, not merely the existence of object files.

\subsection{Recorded checks and limitations}
The baseline source-only verification recorded successful checks for all
181 selected project modules and 23 audit/control steps. Four false Lean
mutations and a certificate coefficient corruption were rejected. The final
axiom reports contain only \code{propext}, \code{Classical.choice}, and
\code{Quot.sound}; there is no \code{sorryAx}, admitted theorem, or numerical
oracle in the public theorem's dependency closure. Kernel-reduced finite
reflection is used explicitly. The original build driver's metadata
finalization was interrupted after the proof steps; a separate bound audit
and finalization are disclosed in \pathcode{verification/BUILD_FINALIZATION.md}.
They are not represented as an observed uninterrupted driver exit.

The pinned environment is Lean 4.34.1 with Mathlib commit
\begin{quote}\small
\pathcode{d13f23b723b8a846827a245b89c10fc7d3f11612}.
\end{quote}
The second implementation, nanoda, accepted all 55,248 exported declarations
in the dependency closure of \code{thomson_nine_complete}, printed the actual
complete target, and permitted only the above three axioms. The checker is
pinned to commit
\pathcode{3a2407216ee84a75f9e1aead6803d0578be06ae7}, and the exporter to
\pathcode{076e8e57707e813375e8f9da8bf989799ace9680}. The recorded run was on
Windows. An earlier timeout remains an inconclusive attempt, not a proof
rejection; the successful retry used the same exported proof. The optional
POSIX recipe has not been executed. Full details and raw output are in
\pathcode{verification/secondary/README.md}.

Independent implementation checking reduces reliance on a single kernel,
but still involves the exporter, implementations, compiler toolchains, and
source-to-export binding. Neither it nor the primary Lean check replaces
human examination of whether the statement captures the intended problem.
No external referee's approval or independent operator's rebuild is claimed.
The paper is a review draft; its explanatory text is not itself kernel checked.

\paragraph{Artifact availability.}
Version 0.1.0 repository:
\url{https://github.com/aravindkarthigeyan/thomson-n9-lean}.
Archive DOI: \href{https://doi.org/10.5281/zenodo.23192843}{10.5281/zenodo.23192843}
(reserved during preparation; deposit and publication are not confirmed by
reservation alone).

\section*{Acknowledgments and assistance disclosure}
This project received substantial assistance from Codex, Claude agents, and
Astra in mathematical exploration, certificate generation, Lean development,
code checking, and preparation of this manuscript. They are tools, not human
coauthors or independent referees. Component reviews by Claude are recorded
in the artifact; Claude stopped before final verification, so no final Claude
approval is implied. The N7 and N8 authors' mathematical and formalization
work was an important foundation for the strategy. Third-party source and
license provenance are retained separately. This manuscript is licensed
CC BY 4.0; original project code uses the license specified in the public
repository, without overriding third-party licenses.

\begin{thebibliography}{9}
\bibitem{cohnkumar}
H. Cohn and A. Kumar, \emph{Universally optimal distribution of points on
spheres}, J. Amer. Math. Soc. \textbf{20} (2007), 99--148.
\href{https://arxiv.org/abs/math/0607446}{arXiv:math/0607446}.

\bibitem{schwartz}
R. E. Schwartz, \emph{The 5 Electron Case of Thomson's Problem},
Experimental Mathematics \textbf{22} (2013).
\href{https://arxiv.org/abs/1001.3702}{arXiv:1001.3702}.

\bibitem{klt}
L. Kryvonos, L. Liehr, and M. A. Taylor,
\emph{Energy minimization for eight points on the sphere} (2026).
\href{https://arxiv.org/abs/2609.22077v1}{arXiv:2609.22077v1}.

\bibitem{n8code}
L. Liehr et al., \emph{Energy-Minimization-8-Points}, source repository.
\url{https://github.com/lukasliehr/Energy-Minimization-8-Points}.
Inspected commit \pathcode{50d14bc06bd41f61573eb6a36eedb8d00759af8f}.
Also inspected: \emph{Thomson-N-8-Warrant},
\url{https://github.com/jstoobysmith/Thomson-N-8-Warrant}, commit
\pathcode{2c35d62e07d72686931da4bc3e062d668e858eda}.

\bibitem{tran}
H. Tran, \emph{thomson-n7-lean}, paper and formal source (2026).
\url{https://github.com/huwngtran/thomson-n7-lean}.
Inspected commit \pathcode{25f2fa53119273458cfbb3c4230904bffdd61e53}.
The proved N7 endpoints and the separately unproved N9 challenge are
distinguished in this source.

\bibitem{delsarte}
\emph{delsarte}, formal spherical-code and Schoenberg positivity library.
\url{https://github.com/specialjcg/delsarte}.
Vendored dependency commit \pathcode{f0724e16b01d4561ee6cda8cd56db4db03de0496};
Apache-2.0, with compatibility changes documented in the artifact.

\bibitem{nanoda}
\emph{nanoda\_lib}, independent Lean proof-term checker.
\url{https://github.com/ammkrn/nanoda_lib};
commit \pathcode{3a2407216ee84a75f9e1aead6803d0578be06ae7}.
Proof export: \url{https://github.com/leanprover/lean4export},
commit \pathcode{076e8e57707e813375e8f9da8bf989799ace9680}.
\end{thebibliography}
\end{document}
